\documentclass[aspectratio=169,12pt]{beamer}
% Maths Lycée — thème éditorial inspiré du portfolio de Yassine.
% Modifier ici les couleurs et la typographie de tous les diaporamas.
\usepackage[T1]{fontenc}
\usepackage[utf8]{inputenc}
\usepackage[french]{babel}
\usepackage{lmodern,amsmath,amssymb,mathtools,tikz,pgfplots,listings,adjustbox}
\usetikzlibrary{arrows.meta,positioning,calc}
\pgfplotsset{compat=1.18,width=9cm,height=4.2cm}
\definecolor{Ink}{HTML}{202124}
\definecolor{Blue}{HTML}{356B9E}
\definecolor{Mist}{HTML}{F6F6F5}
\definecolor{Muted}{HTML}{737579}
\definecolor{Line}{HTML}{E7E7E5}
\setbeamercolor{normal text}{fg=Ink,bg=white}
\setbeamercolor{structure}{fg=Blue}
\setbeamercolor{frametitle}{fg=Ink}
\setbeamercolor{title}{fg=Ink}
\setbeamercolor{block title}{fg=Blue,bg=Mist}
\setbeamercolor{block body}{fg=Ink,bg=Mist}
\setbeamerfont{title}{size=\LARGE,series=\bfseries}
\setbeamerfont{frametitle}{size=\Large,series=\bfseries}
\setbeamersize{text margin left=1cm,text margin right=1cm}
\setbeamertemplate{navigation symbols}{}
\setbeamertemplate{itemize items}[circle]
\setbeamertemplate{enumerate items}[default]
\setlength{\parskip}{7pt}
\newcommand{\slidekind}{COMPRENDRE}
\newcommand{\chapternumber}{01}
\setbeamertemplate{frametitle}{\vspace{0.3cm}{\fontsize{7}{9}\selectfont\color{Blue}\MakeUppercase{\slidekind}}\par\vspace{0.1cm}\insertframetitle\par\vspace{0.08cm}\color{Line}\rule{\textwidth}{0.4pt}\vspace{-0.02cm}}
\setbeamertemplate{footline}{\hspace{1cm}\begin{minipage}{\dimexpr\paperwidth-2cm}\color{Line}\rule{\textwidth}{0.4pt}\par\vspace{0.06cm}{\fontsize{6.5}{8}\selectfont\color{Muted}MATHS LYCÉE\quad / \quad PREMIÈRE SPÉCIALITÉ\hfill\insertframenumber\ /\ \inserttotalframenumber}\end{minipage}\vspace{0.22cm}}
\setbeamertemplate{title page}{
\vspace{0.3cm}
{\fontsize{8}{10}\selectfont\color{Blue}MATHS LYCÉE\quad /\quad 2026--2027}\par
\vspace{0.6cm}
\begin{columns}[c,onlytextwidth]
\begin{column}{.72\textwidth}
{\fontsize{9}{12}\selectfont\color{Muted}CHAPITRE \chapternumber}\par\vspace{0.22cm}
{\usebeamerfont{title}\color{Ink}\inserttitle\par}
\vspace{0.4cm}{\color{Blue}\rule{1.5cm}{2pt}}\par
{\small\color{Muted}Comprendre les notions.\par Justifier chaque étape.}
\end{column}
\begin{column}{.23\textwidth}
\begin{tikzpicture}[scale=.65]
\fill[Mist,rounded corners=6pt](-1.9,-2.1)rectangle(2,2.5);
\draw[Line,step=.5](-1.8,-2)grid(1.9,2.4);
\draw[Muted,->](-1.7,-1.3)--(1.8,-1.3);
\draw[Muted,->](0,-1.6)--(0,2.2);
\draw[Blue,very thick,domain=-1.5:1.5,samples=40]plot(\x,{1.3*\x*\x-1.3});
\fill[Blue](1,0)circle(2.5pt);
\end{tikzpicture}
\end{column}
\end{columns}
\vfill}
\lstdefinestyle{mathspython}{language=Python,basicstyle=\ttfamily\small,keywordstyle=\color{Blue}\bfseries,commentstyle=\color{Muted},stringstyle=\color{Blue},backgroundcolor=\color{Mist},frame=none,xleftmargin=8pt,xrightmargin=8pt,aboveskip=8pt,belowskip=8pt,showstringspaces=false,columns=fullflexible,keepspaces=true,breaklines=true}
\lstnewenvironment{codepython}{\lstset{style=mathspython}}{}
\newcommand{\R}{\mathbb{R}}
\newcommand{\N}{\mathbb{N}}
\newcommand{\E}{\mathbb{E}}
\newcommand{\V}{\mathrm{V}}
\hypersetup{colorlinks=true,urlcolor=Blue}

\title{Variations, extrema et optimisation}
\subtitle{Première spécialité}
\renewcommand{\chapternumber}{06}\author{Maths Lycée}
\date{Année scolaire 2026--2027}
\begin{document}\begin{frame}\titlepage\end{frame}
\begin{frame}{Objectifs du chapitre}\begin{itemize}
\item Relier dérivée et variations sur un intervalle
\item Construire un tableau de variations
\item Déterminer un extremum
\item Modéliser un problème d’optimisation
\item Étudier une parité et démontrer une inégalité
\end{itemize}\end{frame}
\renewcommand{\slidekind}{COMPRENDRE}
\begin{frame}[t,fragile]{Une boîte à optimiser}
On découpe des carrés de côté $x$ dans les coins d’un carton carré de côté $12$ cm, puis on replie les bords. Le volume est $V(x)=x(12-2x)^2$.\medskip\par Les longueurs imposent $0<x<6$. Augmenter la hauteur diminue la base : le meilleur choix ne se lit pas directement dans la formule.
\end{frame}
\renewcommand{\slidekind}{COMPRENDRE}
\begin{frame}[t,fragile]{Croissance sur un intervalle}
$f$ est croissante sur $I$ si, pour tous $a<b$ dans $I$, on a $f(a)\le f(b)$. Elle est strictement croissante si $f(a)<f(b)$.\medskip\par Le mot « sur » est essentiel : une fonction peut croître sur un intervalle et décroître sur un autre.
\end{frame}
\renewcommand{\slidekind}{COMPRENDRE}
\begin{frame}[t,fragile]{Le théorème des variations}
Soit $f$ dérivable sur un intervalle $I$.\begin{itemize}\item Si $f\prime\ge0$ sur $I$, alors $f$ est croissante sur $I$.\item Si $f\prime\le0$, elle est décroissante.\item Si $f\prime=0$ partout, elle est constante.\end{itemize} Ce théorème est admis en première. Une dérivée strictement positive assure une croissance stricte.
\end{frame}
\renewcommand{\slidekind}{COMPRENDRE}
\begin{frame}[t,fragile]{Une condition portant sur un intervalle}
$f\prime(2)>0$ décrit une pente en un point. Cela ne suffit pas à conclure que $f$ croît sur tout son domaine.\medskip\par Pour $f(x)=x^2$, $f\prime(2)=4$, mais $f$ décroît sur $]-\infty,0]$. On doit étudier le signe de $f\prime(x)$ pour toutes les valeurs de l’intervalle.
\end{frame}
\renewcommand{\slidekind}{COMPRENDRE}
\begin{frame}[t,fragile]{Étudier le signe avant de dessiner les flèches}
Pour $f(x)=x^3-3x$,\[f\prime(x)=3x^2-3=3(x-1)(x+1).\] Les zéros sont $-1$ et $1$. La dérivée est positive à l’extérieur de $[-1,1]$ et négative à l’intérieur. Les flèches sont donc : montée, descente, montée.
\end{frame}
\renewcommand{\slidekind}{COMPRENDRE}
\begin{frame}[t,fragile]{Un tableau complet sur un segment}
Pour $f(x)=x^3-3x$ sur $[-2,2]$,\[\begin{array}{c|cccc}x&-2&-1&1&2\\\hline f(x)&-2&2&-2&2\end{array}\]\medskip Les variations successives sont $-2\nearrow2$, puis $2\searrow-2$, puis $-2\nearrow2$. Le tableau doit présenter les valeurs aux bornes et aux points critiques, pas seulement le signe de la dérivée.
\end{frame}
\renewcommand{\slidekind}{COMPRENDRE}
\begin{frame}[t,fragile]{Extremum local et global}
Un maximum local est une valeur supérieure aux valeurs voisines. Un maximum global domine toutes les valeurs du domaine étudié.\medskip\par Pour $x^3-3x$ sur $\R$, $f(-1)=2$ est un maximum local, mais il n’est pas global : $f(3)=18>2$. Le domaine change la conclusion.
\end{frame}
\renewcommand{\slidekind}{COMPRENDRE}
\begin{frame}[t,fragile]{Un changement de signe décisif}
Quand $f\prime$ passe de $+$ à $-$ en $a$, la fonction monte puis descend : $f(a)$ est un maximum local. Le passage de $-$ à $+$ donne un minimum local.\medskip\par Si $f\prime(a)=0$ sans changement de signe, aucun extremum n’est garanti.
\end{frame}
\renewcommand{\slidekind}{COMPRENDRE}
\begin{frame}[t,fragile]{Une dérivée nulle sans extremum}
$f(x)=x^3$ vérifie $f\prime(x)=3x^2\ge0$. La dérivée s’annule en $0$, mais reste positive de chaque côté. La fonction continue à croître.\medskip\par Ainsi « $f\prime(a)=0$ » est une piste pour chercher un extremum intérieur, pas une conclusion.
\end{frame}
\renewcommand{\slidekind}{COMPRENDRE}
\begin{frame}[t,fragile]{Les bornes comptent aussi}
Pour $f(x)=x^2$ sur $[-1,3]$, le point critique est $0$. On compare\[f(-1)=1,\qquad f(0)=0,\qquad f(3)=9.\] Le minimum vaut $0$ et le maximum $9$. Le maximum est atteint au bord, où la dérivée n’est pas nulle.
\end{frame}
\renewcommand{\slidekind}{MÉTHODE \& EXEMPLE}
\begin{frame}[t,fragile]{La méthode d’optimisation}
\begin{enumerate}\item Choisir une variable et préciser les contraintes.\item Exprimer la quantité à optimiser en fonction de cette variable.\item Dériver, factoriser et étudier le signe.\item Comparer les valeurs candidates, y compris les bornes.\item Répondre dans le contexte, avec les unités.\end{enumerate}
\end{frame}
\renewcommand{\slidekind}{MÉTHODE \& EXEMPLE}
\begin{frame}[t,fragile]{La boîte : calcul de la dérivée}
$V(x)=x(12-2x)^2$ sur $]0,6[$. Par produit et composition affine,\[V\prime(x)=(12-2x)^2-4x(12-2x)=(12-2x)(12-6x).\] Garder la factorisation rend le signe immédiat : le premier facteur est positif dans le domaine.
\end{frame}
\renewcommand{\slidekind}{COMPRENDRE}
\begin{frame}[t,fragile]{La boîte : décision et unités}
$12-6x$ est positif pour $x<2$ et négatif pour $x>2$. Le volume augmente jusqu’à $2$ puis diminue.\[V(2)=2\times8^2=128\ \text{cm}^3.\] La boîte optimale mesure $8$ cm sur $8$ cm, avec une hauteur de $2$ cm. La valeur $128$ est un volume, pas une longueur.
\end{frame}
\renewcommand{\slidekind}{COMPRENDRE}
\begin{frame}[t,fragile]{Fonctions paires}
Le domaine doit être symétrique par rapport à zéro. Une fonction est paire si $f(-x)=f(x)$ pour tout $x$ du domaine.\medskip\par La courbe possède une symétrie par rapport à l’axe des ordonnées. Par exemple $x^4-2x^2$ est paire. On peut étudier les variations sur $[0,+\infty[$ puis utiliser cette symétrie.
\end{frame}
\renewcommand{\slidekind}{COMPRENDRE}
\begin{frame}[t,fragile]{Fonctions impaires}
Sur un domaine symétrique, $f$ est impaire si $f(-x)=-f(x)$. Sa courbe possède une symétrie centrale de centre l’origine.\medskip\par $x^3-3x$ est impaire. Une fonction n’est pas nécessairement paire ou impaire : $x^2+x$ n’est ni l’une ni l’autre.
\end{frame}
\renewcommand{\slidekind}{COMPRENDRE}
\begin{frame}[t,fragile]{Comparer deux fonctions}
Pour montrer $f(x)\ge g(x)$ sur $I$, étudier $h=f-g$ et chercher son minimum.\medskip\par Exemple : $x^2\ge2x-1$ sur $\R$. Ici $h=x^2-2x+1$ et $h\prime=2x-2$. Le minimum est atteint en $1$ et vaut $0$. Donc $h\ge0$, ce qui prouve l’inégalité.
\end{frame}
\renewcommand{\slidekind}{COMPRENDRE}
\begin{frame}[t,fragile]{Un résultat exact avant l’arrondi}
Un problème peut donner une valeur optimale irrationnelle. Conserver cette valeur pendant les calculs, puis arrondir seulement la réponse finale.\medskip\par Si la variable représente un nombre entier d’objets, comparer les entiers admissibles proches de l’optimum continu. On ne peut pas produire $3,7$ objets.
\end{frame}
\renewcommand{\slidekind}{À RETENIR}
\begin{frame}[t,fragile]{Auto-vérification}
Un tableau cohérent donne des valeurs plus grandes au sommet d’une flèche montante. Un extremum doit respecter les contraintes du problème.\medskip\par Avant de rendre : ai-je précisé le domaine ? justifié le signe ? examiné les bornes ? distingué abscisse optimale et valeur optimale ? répondu avec une phrase et une unité ?
\end{frame}
\renewcommand{\slidekind}{COMPRENDRE}
\begin{frame}[t,fragile]{Lire un maximum sur la courbe}
\begin{columns}[c,onlytextwidth]\begin{column}{.60\textwidth}\centering\begin{adjustbox}{max width=\linewidth,max height=3.7cm}\begin{tikzpicture}\begin{axis}[width=9cm,height=4.4cm,axis lines=left,xmin=0,xmax=6,ymin=0,ymax=140,xlabel={$x$ (cm)},ylabel={$V$ (cm$^3$)}]\addplot[Blue,thick,domain=0:6,samples=70]{x*(12-2*x)^2};\addplot[Muted,dashed]coordinates{(2,0)(2,128)(0,128)};\addplot[only marks]coordinates{(2,128)};\end{axis}\end{tikzpicture}\end{adjustbox}\end{column}\begin{column}{.36\textwidth} Le sommet confirme le calcul : hauteur $2$ cm, volume $128$ cm$^3$. Le graphique aide à interpréter ; le signe de la dérivée justifie le maximum.\end{column}\end{columns}
\end{frame}
\renewcommand{\slidekind}{À VOUS DE JOUER}\begin{frame}[t]{Vérification des acquis 1}
$f\prime$ passe de $-$ à $+$ en $2$. On obtient :\begin{enumerate}
\item Un maximum
\item Un minimum
\item Aucun extremum
\item Une racine de $f$
\end{enumerate}\medskip Une seule réponse est correcte. Justifier votre choix.\end{frame}
\renewcommand{\slidekind}{COMPRENDRE LA RÉPONSE}\begin{frame}[t]{Vérification 1 : résolution}
\begin{block}{Réponse 2}Un minimum\end{block}\medskip La fonction descend puis monte : elle atteint un minimum local.\end{frame}
\renewcommand{\slidekind}{À VOUS DE JOUER}\begin{frame}[t]{Vérification des acquis 2}
Pour chercher un maximum sur $[a,b]$, il faut aussi examiner :\begin{enumerate}
\item Seulement $0$
\item Seulement $f\prime(a)$
\item Aucun autre point
\item Les bornes $a$ et $b$
\end{enumerate}\medskip Une seule réponse est correcte. Justifier votre choix.\end{frame}
\renewcommand{\slidekind}{COMPRENDRE LA RÉPONSE}\begin{frame}[t]{Vérification 2 : résolution}
\begin{block}{Réponse 4}Les bornes $a$ et $b$\end{block}\medskip Un extremum global peut être atteint à une borne.\end{frame}
\end{document}
